% from main.tex
%%%%%%%%%%  scan idx 286  |  folio 273  %%%%%%%%%%
peut développer pour les espaces analytiques complexes un tapis analogue à celui de ce numéro (propriétés d'échange, de semi-continuité, de continuité)\footnote{C'est d'ailleurs dans ce cadre que ces propriétés ont été d'abord formulées (voir par exemple Grauert [1]), avant d'être transposées en Géométrie Algébrique dans (EGA III 7).}. Les ingrédients essentiels sont : 1) le théorème de finitude pour les morphismes propres («~théorème de Grauert~») ; 2) un analogue de la formule de projection (3.7), qu'on peut énoncer ainsi : soit \(f:X\to Y\) un morphisme d'espaces analytiques, où \(X\) est supposé de dimension majorée, et soient \(E\in\Ob D^-(X)\), \(M\in\Ob D^-(Y)_{\coh}\) («~coh~» signifie «~à cohomologie cohérente~»). La flèche canonique
\[
 Rf_*(E)\otimes^{\mathrm L}M)\longrightarrow Rf_*(E\otimes_Y^{\mathrm L}M)
\]
est un isomorphime (l'hypothèse sur \(X\) assure (T.F II 5.12) que \(f_*\) est de dimension cohomologique finie, donc que \(Rf_*\) est défini sur \(D(X)\)). A la différence de l'énoncé 1), l'énoncé 2) est essentiellement trivial, et résulte d'un dévissage analogue à celui de la preuve de (3.7) ; on a d'ailleurs les coudées plus franches puisque \(Rf_*\) est défini sur \(D(X)\), en particulier il n'est besoin d'aucune hypothèse sur \(E\) (en dehors de \(E\in\Ob D^-(X)\)).
\end{enumerate}

\subsection*{Bibliographie}
\begin{enumerate}[label={[\arabic*]}]
\item[{[1]}] Grauert, H., \emph{Ein Theorem der Analytischen Garbentheorie} \ldots, Pub. IHES, n\textsuperscript{o}5.
\item[{[H]}] Hartshorne, R., \emph{Residues and Duality}, Lecture Notes n\textsuperscript{o}20, Springer, 1966.
\item[{[T.F]}] Godement, R., \emph{Théorie des faisceaux}, Hermann, 1958.
\end{enumerate}



% Exposé IV macros
\renewcommand{\Ob}{\operatorname{ob}}
\providecommand{\Mod}{\operatorname{Mod}}
\providecommand{\Coh}{\operatorname{Coh}}
\providecommand{\Loclib}{\operatorname{Loclib}}
\providecommand{\coh}{\mathrm{coh}}
\providecommand{\parf}{\mathrm{parf}}
\providecommand{\naif}{\text{naïf}}
\providecommand{\loc}{\mathrm{loc}}
\providecommand{\fin}{\mathrm{fin}}
\providecommand{\clA}{\operatorname{cl}}
\providecommand{\Rg}{\operatorname{rg}}
\providecommand{\varinjlim}{\varinjlim}
\providecommand{\varprojlim}{\varprojlim}
\renewcommand{\id}{\operatorname{id}}
\providecommand{\Th}{\operatorname{Th}}
\providecommand{\Hom}{\operatorname{Hom}}
\providecommand{\Dcat}{\mathrm D}
\providecommand{\Kcat}{\mathrm K}
\providecommand{\Kdot}{K_{\bullet}}
\providecommand{\Kprime}{K'}
\providecommand{\naif}{\text{naïf}}


